%% This is file `elsarticle-template-1-num.tex',
%%
%% Copyright 2009 Elsevier Ltd
%%
%% This file is part of the 'Elsarticle Bundle'.
%% ---------------------------------------------
%%
%% It may be distributed under the conditions of the LaTeX Project Public
%% License, either version 1.2 of this license or (at your option) any
%% later version.  The latest version of this license is in
%%    http://www.latex-project.org/lppl.txt
%% and version 1.2 or later is part of all distributions of LaTeX
%% version 1999/12/01 or later.
%%
%% The list of all files belonging to the 'Elsarticle Bundle' is
%% given in the file `manifest.txt'.
%%
%% Template article for Elsevier's document class `elsarticle'
%% with numbered style bibliographic references
%%

\documentclass[1p,12pt]{elsarticle}

\usepackage{amssymb}
\usepackage{lineno}
\usepackage{epsfig}
\usepackage[spanish]{babel}
\usepackage[labelsep=endash]{caption}
\usepackage{float}
\usepackage{booktabs}
\usepackage{multirow}
\usepackage{subcaption}

%% Paquetes para evitar desbordamiento
\usepackage{microtype}       % mejora el espaciado tipográfico y reduce overfull hbox
\usepackage{array}           % columnas con ancho fijo en tablas
\usepackage{tabularx}        % tablas que se ajustan al ancho del texto
\usepackage{adjustbox}       % redimensionar tablas que se salen del margen
\usepackage[hyphens]{url}    % permite partir URLs en saltos de línea

%% Permite que LaTeX parta palabras largas en código monoespaciado
\usepackage{listings}
\lstset{breaklines=true, breakatwhitespace=true, basicstyle=\small\ttfamily}

%% Evita que hbox overful queden sin advertencia silenciosa
\setlength{\emergencystretch}{3em}

\journal{Cátedra de Inteligencia Artificial}

\renewcommand{\figurename}{Fig.}
\renewcommand{\tablename}{Tab.}

\begin{document}

\begin{frontmatter}

\title{Clasificación de Imágenes de Animales mediante Redes Neuronales
       Convolucionales \\ Trabajo Práctico N°3}

\autor{Arnedo Emmanuel, García Iñaki, Hamada Santino, Veliz Matías}

\catedra{Cátedra de Inteligencia Artificial \\ Ingeniería en Sistemas de
         Información}

\begin{abstract}
Se presenta la comparación de tres arquitecturas de redes neuronales
convolucionales (CNN) para la clasificación de imágenes de animales en diez
categorías: araña, ardilla, caballo, elefante, gallina, gato, mariposa, oveja,
perro y vaca. Se implementaron un modelo Custom CNN diseñado desde cero y dos
modelos basados en Transfer Learning utilizando VGG16~\cite{Simonyan2015} y
MobileNetV2~\cite{Sandler2018} preentrenados en ImageNet~\cite{Deng2009}. La
evaluación se realizó sobre el conjunto de test mediante métricas de
clasificación multiclase: Accuracy, Precision, Recall, F1-Score, Specificity y
ROC-AUC (One-vs-Rest). Los tres modelos alcanzaron desempeños variados, siendo
VGG16 el de mejor desempeño con un Accuracy de 0.9612 y el menor costo
computacional, lo que lo posiciona como la alternativa más eficiente para este
problema.
\end{abstract}

\begin{keyword}
Redes Neuronales Convolucionales \sep Transfer Learning \sep Clasificación de
Imágenes \sep VGG16 \sep MobileNetV2 \sep Clasificación Multiclase
\end{keyword}

\end{frontmatter}

%% main text
\section{Introducción}
\label{sec:intro}

El presente trabajo aborda el problema de clasificación de imágenes de animales
en diez especies distintas mediante redes neuronales convolucionales
(CNN)~\cite{LeCun2015}. El objetivo es comparar tres arquitecturas diferentes y
evaluar su rendimiento utilizando métricas robustas para clasificación
multiclase, evitando depender exclusivamente de la métrica Accuracy.

Se implementaron tres modelos con estrategias diferenciadas:

\begin{enumerate}
    \item \textbf{Custom CNN:} arquitectura convolucional diseñada e implementada
          desde cero, con cuatro bloques convolucionales y una cabeza de
          clasificación densa.
    \item \textbf{VGG16:} modelo preentrenado en ImageNet~\cite{Deng2009}
          utilizado como extractor de características (Feature Extraction), con
          todas las capas de la base congeladas~\cite{Simonyan2015}.
    \item \textbf{MobileNetV2:} modelo preentrenado en ImageNet diseñado para
          eficiencia computacional, también utilizado en modo Feature
          Extraction~\cite{Sandler2018}.
\end{enumerate}

La técnica de Transfer Learning~\cite{Pan2010} permite reutilizar el
conocimiento aprendido por modelos entrenados en conjuntos de datos masivos
(ImageNet: más de 14 millones de imágenes en 1000 clases) y aplicarlo a
problemas más específicos y con menor disponibilidad de datos.

La evaluación de los modelos se realizó sobre un conjunto de test independiente,
empleando un conjunto de métricas que incluye Accuracy, Precision (macro),
Recall (macro), F1-Score (macro), Specificity (macro) y ROC-AUC
(One-vs-Rest)~\cite{Sokolova2009}. Se priorizó el F1-Score Macro como métrica
principal debido a que permite una evaluación equilibrada del rendimiento por
clase, lo cual resulta relevante ante posibles desbalances en la distribución de
las clases.

\section{Dataset}
\label{sec:dataset}

Se utilizó el dataset \textbf{Animals Dataset}, compuesto por 26.179 imágenes
organizadas en 10 clases desbalanceadas: araña, ardilla, caballo, elefante,
gallina, gato, mariposa, oveja, perro y vaca.

Todas las imágenes se redimensionaron a $224 \times 224$ píxeles para
garantizar la compatibilidad con las tres arquitecturas evaluadas.

\subsection{Distribución del dataset}
\label{sub:distribucion}

El dataset se dividió en tres subconjuntos, cuya distribución se presenta en la
Tab.~\ref{tab:distribucion_dataset}.

\begin{table}[h]
\centering
\begin{tabular}{lrr}
\hline
\textbf{Subconjunto} & \textbf{Total de imágenes} &
\textbf{Imágenes por clase} \\
\hline
Train      & 20938 & Desbalanceado \\
Validation & 2614  & Desbalanceado \\
Test       & 2627  & Desbalanceado \\
\hline
\end{tabular}
\caption{Distribución del dataset Animals Dataset Splitted.}
\label{tab:distribucion_dataset}
\end{table}

\subsection{Data Augmentation}
\label{sub:data_augmentation}

Para mejorar la capacidad de generalización de los modelos y reducir el
overfitting, se aplicaron transformaciones aleatorias de Data
Augmentation~\cite{Shorten2019} \textbf{únicamente al conjunto de
entrenamiento}. Los parámetros configurados se detallan en la
Tab.~\ref{tab:data_augmentation}.

% Tabla con adjustbox para evitar desbordamiento
\begin{table}[h]
\centering
\begin{adjustbox}{max width=\textwidth}
\begin{tabular}{lll}
\hline
\textbf{Parámetro} & \textbf{Valor} & \textbf{Descripción} \\
\hline
\texttt{rotation\_range}      & 20°  & Rotación aleatoria de hasta $\pm$20 grados \\
\texttt{width\_shift\_range}  & 0.2  & Desplazamiento horizontal de hasta 20\% \\
\texttt{height\_shift\_range} & 0.2  & Desplazamiento vertical de hasta 20\% \\
\texttt{zoom\_range}          & 0.2  & Zoom aleatorio de hasta 20\% \\
\texttt{horizontal\_flip}     & True & Espejo horizontal aleatorio \\
\hline
\end{tabular}
\end{adjustbox}
\caption{Parámetros de Data Augmentation aplicados al conjunto de entrenamiento.}
\label{tab:data_augmentation}
\end{table}

\section{Propuesta de Desarrollo}
\label{sec:propuesta}

Se propuso implementar tres arquitecturas CNN con estrategias diferenciadas para
abordar el problema de clasificación multiclase de imágenes de animales: una red
diseñada desde cero (Custom CNN) y dos redes preentrenadas en ImageNet
aplicando Transfer Learning en modo Feature Extraction (VGG16 y MobileNetV2).

Todos los modelos comparten la siguiente configuración base de entrenamiento:

\begin{itemize}
    \item \textbf{Loss:} Sparse Categorical Crossentropy.
    \item \textbf{Métrica de monitoreo:} Accuracy.
    \item \textbf{Epochs máximos:} 50.
    \item \textbf{EarlyStopping:} \texttt{patience=5}, monitorea
          \texttt{val\_accuracy},\\ \texttt{restore\_best\_weights=True}.
    \item \textbf{Batch size:} 32.
\end{itemize}

\subsection{Custom CNN}
\label{sub:custom_cnn}

Red convolucional diseñada e implementada desde cero con la siguiente
arquitectura:

\begin{itemize}
    \item \texttt{Conv2D(32,\,3,\,``same'',\,ReLU)} $\rightarrow$
          \texttt{BatchNorm()} $\rightarrow$ \texttt{MaxPool2D()}.
    \item \texttt{Conv2D(64,\,3,\,``same'',\,ReLU)} $\rightarrow$
          \texttt{BatchNorm()} $\rightarrow$ \texttt{MaxPool2D()}.
    \item \texttt{Conv2D(128,\,3,\,``same'',\,ReLU)} $\rightarrow$
          \texttt{BatchNorm()} $\rightarrow$ \texttt{MaxPool2D()}.
    \item \texttt{Conv2D(256,\,3,\,``same'',\,ReLU)} $\rightarrow$
          \texttt{BatchNorm()}.
    \item \texttt{GlobalAveragePooling2D()}.
    \item \texttt{Dense(256,\,ReLU)} $\rightarrow$ \texttt{BatchNorm()}
          $\rightarrow$ \texttt{Dropout(0.3)}.
    \item \texttt{Dense(10,\,Softmax)}: capa de salida con 10 clases.
\end{itemize}

El optimizador utilizado fue Adam~\cite{Kingma2015} con un learning rate
explícito de $lr = 1 \times 10^{-4}$ para asegurar una convergencia más
estable. La normalización de las imágenes mediante
\texttt{rescale=1./255} se delegó al generador de datos
(\texttt{ImageDataGenerator}) durante el entrenamiento. Esto implica que,
durante la inferencia, es necesario aplicar este preprocesamiento de manera
externa (dividiendo los píxeles por 255) antes de inyectar los tensores a la
red.

\subsection{Manejo del Desbalanceo: Class Weights}
\label{sub:class_weights}

Dado que el \textit{Animals Dataset} presenta un fuerte desbalanceo de clases
(por ejemplo, la clase \textit{perro} posee 4863 imágenes mientras que
\textit{elefante} cuenta con solo 1446), entrenar a los modelos sin ninguna
compensación provocaría un sesgo hacia las clases mayoritarias.

Para mitigar este problema sin recurrir a técnicas de submuestreo
(undersampling) o sobremuestreo (oversampling) que alteren el dataset original,
se utilizó el parámetro \texttt{class\_weight} proporcionado por Keras durante
el entrenamiento. Los pesos de las clases se calcularon matemáticamente mediante
la heurística de balanceo de scikit-learn (\texttt{compute\_class\_weight}),
cuya fórmula es inversamente proporcional a la frecuencia de cada clase:

\begin{equation}
\label{eq:class_weight}
w_j = \frac{N}{K \cdot n_j}
\end{equation}

donde $w_j$ es el peso asignado a la clase $j$, $N$ es el total de muestras del
dataset, $K$ es el número total de clases, y $n_j$ es la cantidad de muestras
pertenecientes a la clase $j$. Al pasar este diccionario de pesos a la red
neuronal, la función de pérdida penaliza de forma mucho más severa los errores
cometidos en clases minoritarias, obligando al modelo a prestarles igual
atención durante el descenso de gradiente.

\subsection{VGG16 con Transfer Learning}
\label{sub:vgg16}

Se utilizó la red VGG16~\cite{Simonyan2015}, preentrenada en ImageNet,
aplicando Transfer Learning en modo Feature Extraction. Todas las capas de la
base fueron congeladas (\texttt{layer.trainable = False}), cargando el modelo
con \texttt{include\_top=False} para excluir las capas densas originales de
clasificación.

La base congelada de VGG16 cuenta con aproximadamente 14.7 millones de
parámetros no entrenables, organizados en 5 bloques de capas convolucionales
con filtros de tamaño $3 \times 3$ y operaciones de MaxPooling2D.

La cabeza de clasificación añadida consiste en:

\begin{itemize}
    \item \texttt{GlobalAveragePooling2D()}.
    \item \texttt{Dense(256,\,ReLU)} $\rightarrow$ \texttt{Dropout(0.3)}.
    \item \texttt{Dense(10,\,Softmax)}: capa de salida con 10 clases.
\end{itemize}

El optimizador utilizado fue Adam con un learning rate reducido
($lr = 1 \times 10^{-4}$), necesario para no destruir las representaciones
preaprendidas. El preprocesamiento se realizó mediante la función
\texttt{preprocess\_input} de VGG16, que centra los píxeles en cero utilizando
la media de ImageNet por canal RGB.

\subsection{MobileNetV2 con Transfer Learning}
\label{sub:mobilenetv2}

Se utilizó MobileNetV2~\cite{Sandler2018}, arquitectura diseñada para operar
eficientemente en dispositivos con recursos limitados. Fue preentrenada en
ImageNet y aplicada en modo Feature Extraction, con todas las capas de la base
congeladas.

La innovación clave de MobileNetV2 radica en el uso de \textbf{Depthwise
Separable Convolutions}, que separan la convolución espacial de la combinación
de canales reduciendo drásticamente el costo computacional, y \textbf{Inverted
Residual Blocks} con conexiones residuales~\cite{He2016}. La base cuenta con
aproximadamente 3.4 millones de parámetros.

La cabeza de clasificación añadida consiste en:

\begin{itemize}
    \item \texttt{GlobalAveragePooling2D()}.
    \item \texttt{Dense(128,\,ReLU)} $\rightarrow$ \texttt{Dropout(0.3)}.
    \item \texttt{Dense(10,\,Softmax)}: capa de salida con 10 clases.
\end{itemize}

El optimizador utilizado fue Adam con $lr = 1 \times 10^{-4}$. El
preprocesamiento se realizó mediante la función \texttt{preprocess\_input} de
MobileNetV2, que normaliza los píxeles al rango $[-1, 1]$.

\section{Experimentos y Resultados}
\label{sec:experimentos}

Los tres modelos fueron entrenados con las configuraciones descritas en la
Sección~\ref{sec:propuesta} y evaluados sobre el conjunto de test (2627
imágenes, desbalanceadas). La evaluación se realizó utilizando las siguientes
métricas de clasificación multiclase: Accuracy, Precision Macro, Recall Macro,
F1-Score Macro, Specificity Macro y ROC-AUC One-vs-Rest.

\subsection{Curvas de entrenamiento}
\label{sub:curvas_entrenamiento}

Las curvas de Accuracy y Loss durante el entrenamiento y la validación permiten
analizar el comportamiento de cada modelo a lo largo de las epochs,
identificando convergencia y posibles indicios de sobreajuste.

\begin{figure}[H]
\centering
\begin{subfigure}[b]{0.48\textwidth}
\centering
\includegraphics[width=\textwidth]{imagenes/training_custom_acc.png}
\caption{Accuracy en train y validation.}
\label{fig:training_custom_acc}
\end{subfigure}
\hfill
\begin{subfigure}[b]{0.48\textwidth}
\centering
\includegraphics[width=\textwidth]{imagenes/training_custom_loss.png}
\caption{Loss en train y validation.}
\label{fig:training_custom_loss}
\end{subfigure}
\caption{Curvas de entrenamiento del modelo Custom CNN.}
\label{fig:training_custom}
\end{figure}

\begin{figure}[H]
\centering
\begin{subfigure}[b]{0.48\textwidth}
\centering
\includegraphics[width=\textwidth]{imagenes/training_vgg16_acc.png}
\caption{Accuracy en train y validation.}
\label{fig:training_vgg16_acc}
\end{subfigure}
\hfill
\begin{subfigure}[b]{0.48\textwidth}
\centering
\includegraphics[width=\textwidth]{imagenes/training_vgg16_loss.png}
\caption{Loss en train y validation.}
\label{fig:training_vgg16_loss}
\end{subfigure}
\caption{Curvas de entrenamiento del modelo VGG16.}
\label{fig:training_vgg16}
\end{figure}

\begin{figure}[H]
\centering
\begin{subfigure}[b]{0.48\textwidth}
\centering
\includegraphics[width=\textwidth]{imagenes/training_mobilenet_acc.png}
\caption{Accuracy en train y validation.}
\label{fig:training_mobilenet_acc}
\end{subfigure}
\hfill
\begin{subfigure}[b]{0.48\textwidth}
\centering
\includegraphics[width=\textwidth]{imagenes/training_mobilenet_loss.png}
\caption{Loss en train y validation.}
\label{fig:training_mobilenet_loss}
\end{subfigure}
\caption{Curvas de entrenamiento del modelo MobileNetV2.}
\label{fig:training_mobilenet}
\end{figure}

\subsection{Resultados de Accuracy y Loss en Test}
\label{sub:resultados_accuracy}

La Tab.~\ref{tab:resultados_test} presenta los resultados de Accuracy y Loss
obtenidos por cada modelo sobre el conjunto de test.

\begin{table}[h]
\centering
\begin{tabular}{lcc}
\hline
\textbf{Modelo} & \textbf{Accuracy} & \textbf{Loss} \\
\hline
VGG16       & \textbf{0.9612} & 0.1417 \\
MobileNetV2 & 0.9558          & 0.1414 \\
Custom CNN  & 0.7168          & 0.8774 \\
\hline
\end{tabular}
\caption{Resultados de Accuracy y Loss en el conjunto de test.}
\label{tab:resultados_test}
\end{table}

\subsection{Matrices de confusión}
\label{sub:matrices_confusion}

Las matrices de confusión permiten identificar las clases que cada modelo
confunde con mayor frecuencia, proporcionando una visión detallada del
desempeño por categoría. Las Fig.~\ref{fig:cm_custom},
Fig.~\ref{fig:cm_vgg16} y Fig.~\ref{fig:cm_mobilenet} presentan las matrices
obtenidas para cada modelo sobre el conjunto de test.

\begin{figure}[H]
\centering
\includegraphics[width=0.63\textwidth]{imagenes/matriz_confusion_custom.png}
\caption{Matriz de confusión del modelo Custom CNN.}
\label{fig:cm_custom}
\end{figure}

\begin{figure}[H]
\centering
\includegraphics[width=0.63\textwidth]{imagenes/matriz_confusion_vgg16.png}
\caption{Matriz de confusión del modelo VGG16.}
\label{fig:cm_vgg16}
\end{figure}

\begin{figure}[H]
\centering
\includegraphics[width=0.63\textwidth]{imagenes/matriz_confusion_mobilenetv2.png}
\caption{Matriz de confusión del modelo MobileNetV2.}
\label{fig:cm_mobilenet}
\end{figure}
\subsection{Classification Report}
\label{sub:classification_report}

Se generó el Classification Report completo para cada modelo utilizando
\texttt{sklearn.metrics.classification\_report}. A continuación se presentan
las tablas de reporte por clase para cada arquitectura.

\begin{table}[h]
\centering
\begin{adjustbox}{max width=\textwidth}
\begin{tabular}{lcccc}
\hline
\textbf{Clase} & \textbf{Precision} & \textbf{Recall} & \textbf{F1-Score} &
\textbf{Support} \\
\hline
Araña       & 0.90 & 0.78 & 0.83 & 483 \\
Ardilla     & 0.68 & 0.73 & 0.71 & 187 \\
Caballo     & 0.86 & 0.63 & 0.73 & 263 \\
Elefante    & 0.41 & 0.95 & 0.57 & 146 \\
Gallina     & 0.81 & 0.83 & 0.82 & 311 \\
Gato        & 0.58 & 0.61 & 0.59 & 168 \\
Mariposa    & 0.73 & 0.91 & 0.81 & 212 \\
Oveja       & 0.63 & 0.77 & 0.70 & 182 \\
Perro       & 0.87 & 0.48 & 0.62 & 487 \\
Vaca        & 0.63 & 0.73 & 0.67 & 188 \\
\hline
\textbf{Macro avg}    & 0.71 & 0.74 & 0.70 & 2627 \\
\textbf{Weighted avg} & 0.76 & 0.72 & 0.72 & 2627 \\
\hline
\end{tabular}
\end{adjustbox}
\caption{Classification Report del modelo Custom CNN.}
\label{tab:report_custom_cnn}
\end{table}

\begin{table}[h]
\centering
\begin{adjustbox}{max width=\textwidth}
\begin{tabular}{lcccc}
\hline
\textbf{Clase} & \textbf{Precision} & \textbf{Recall} & \textbf{F1-Score} &
\textbf{Support} \\
\hline
Araña       & 1.00 & 0.97 & 0.98 & 483 \\
Ardilla     & 0.96 & 0.96 & 0.96 & 187 \\
Caballo     & 0.95 & 0.93 & 0.94 & 263 \\
Elefante    & 0.95 & 0.99 & 0.97 & 146 \\
Gallina     & 0.99 & 0.99 & 0.99 & 311 \\
Gato        & 0.95 & 0.96 & 0.96 & 168 \\
Mariposa    & 0.94 & 1.00 & 0.97 & 212 \\
Oveja       & 0.89 & 0.97 & 0.93 & 182 \\
Perro       & 0.98 & 0.94 & 0.96 & 487 \\
Vaca        & 0.90 & 0.92 & 0.91 & 188 \\
\hline
\textbf{Macro avg}    & 0.95 & 0.96 & 0.96 & 2627 \\
\textbf{Weighted avg} & 0.96 & 0.96 & 0.96 & 2627 \\
\hline
\end{tabular}
\end{adjustbox}
\caption{Classification Report del modelo VGG16.}
\label{tab:report_vgg16}
\end{table}

\begin{table}[h]
\centering
\begin{adjustbox}{max width=\textwidth}
\begin{tabular}{lcccc}
\hline
\textbf{Clase} & \textbf{Precision} & \textbf{Recall} & \textbf{F1-Score} &
\textbf{Support} \\
\hline
Araña       & 1.00 & 0.98 & 0.99 & 483 \\
Ardilla     & 0.97 & 0.94 & 0.96 & 187 \\
Caballo     & 0.96 & 0.92 & 0.94 & 263 \\
Elefante    & 0.93 & 0.96 & 0.95 & 146 \\
Gallina     & 0.97 & 0.99 & 0.98 & 311 \\
Gato        & 0.95 & 0.97 & 0.96 & 168 \\
Mariposa    & 0.97 & 0.99 & 0.98 & 212 \\
Oveja       & 0.86 & 0.93 & 0.89 & 182 \\
Perro       & 0.98 & 0.94 & 0.96 & 487 \\
Vaca        & 0.87 & 0.90 & 0.88 & 188 \\
\hline
\textbf{Macro avg}    & 0.95 & 0.95 & 0.95 & 2627 \\
\textbf{Weighted avg} & 0.96 & 0.96 & 0.96 & 2627 \\
\hline
\end{tabular}
\end{adjustbox}
\caption{Classification Report del modelo MobileNetV2.}
\label{tab:report_mobilenetv2}
\end{table}

\clearpage
\subsection{Curvas ROC}
\label{sub:curvas_roc}

Se calcularon las curvas ROC One-vs-Rest para cada modelo, representando la
Tasa de Verdaderos Positivos frente a la Tasa de Falsos Positivos para cada
clase. Las Fig.~\ref{fig:roc_custom}, Fig.~\ref{fig:roc_vgg16} y
Fig.~\ref{fig:roc_mobilenet} presentan las curvas correspondientes de cada
modelo.

\begin{figure}[H]
\centering
\includegraphics[width=0.63\textwidth]{imagenes/curvas_roc_custom.png}
\caption{Curvas ROC One-vs-Rest del modelo Custom CNN.}
\label{fig:roc_custom}
\end{figure}

\begin{figure}[H]
\centering
\includegraphics[width=0.63\textwidth]{imagenes/curvas_roc_vgg16.png}
\caption{Curvas ROC One-vs-Rest del modelo VGG16.}
\label{fig:roc_vgg16}
\end{figure}

\begin{figure}[H]
\centering
\includegraphics[width=0.63\textwidth]{imagenes/curvas_roc_mobilenetv2.png}
\caption{Curvas ROC One-vs-Rest del modelo MobileNetV2.}
\label{fig:roc_mobilenet}
\end{figure}
\subsection{Tabla comparativa de métricas}
\label{sub:tabla_comparativa}

La Tab.~\ref{tab:comparativa_final} consolida las métricas de clasificación
multiclase obtenidas por cada modelo. Se priorizó el F1-Score Macro como
métrica principal, dado que se trata de un problema de clasificación donde cada
imagen corresponde a una única especie y el F1-Score Macro proporciona una
evaluación equilibrada del rendimiento por clase, siendo robusta ante posibles
desbalances en la distribución.

\begin{table}[h]
\centering
\begin{adjustbox}{max width=\textwidth}
\begin{tabular}{lcccccc}
\hline
\textbf{Modelo} & \textbf{Accuracy} & \textbf{Precision} & \textbf{Recall} &
\textbf{F1 Macro} & \textbf{Specificity} & \textbf{ROC-AUC} \\
 & & \textbf{Macro} & \textbf{Macro} & & & \\
\hline
VGG16       & \textbf{0.9612} & 0.96 & 0.96 & 0.96 & 0.96 & 0.96 \\
MobileNetV2 & 0.9558          & 0.95 & 0.95 & 0.95 & 0.96 & 0.96 \\
Custom CNN  & 0.7168          & 0.71 & 0.74 & 0.70 & 0.76 & 0.72 \\
\hline
\end{tabular}
\end{adjustbox}
\caption{Tabla comparativa de métricas de clasificación multiclase.}
\label{tab:comparativa_final}
\end{table}

\section{Conclusiones}
\label{sec:conclusiones}

A partir de los experimentos realizados, se obtuvieron las siguientes
conclusiones:

\begin{itemize}
    \item Los modelos evaluados demostraron comportamientos muy variados,
          dejando en evidencia la dificultad de trabajar con 10 clases altamente
          desbalanceadas.

    \item \textbf{VGG16} obtuvo el mejor rendimiento, alcanzando un Accuracy de
          0.9612 (96.12\%). Al actuar como un poderoso extractor de
          características gracias a su preentrenamiento masivo en ImageNet,
          logró separar eficazmente las 10 especies animales.

    \item \textbf{MobileNetV2} quedó en un muy cercano segundo lugar con un
          Accuracy de 0.9558 (95.58\%), posicionándose como la alternativa más
          eficiente en el mundo real, dado que requiere una ínfima fracción del
          costo computacional de VGG16 para obtener el mismo resultado práctico.

    \item \textbf{Custom CNN} alcanzó un Accuracy de 0.7168 (71.68\%). A
          pesar de contar con regularización, este modelo entrenado desde cero
          encontró su techo de aprendizaje rápidamente debido a su arquitectura
          más simple y la falta de pesos inicializados. Al enfrentarse a 10
          clases con alta variabilidad intraclase (distintos fondos y poses),
          la red no logró extraer descriptores abstractos tan robustos como las
          redes preentrenadas en ImageNet, estancándose y sin lograr aumentar su
          precisión. Esto demuestra empíricamente la superioridad del
          \textit{Transfer Learning} frente a arquitecturas inicializadas desde
          cero cuando la complejidad de las imágenes aumenta.

    \item La utilización del parámetro \texttt{class\_weight} en el
          entrenamiento fue de vital importancia. Sin esta penalización a nivel
          función de pérdida, los modelos habrían colapsado estadísticamente
          hacia las clases mayoritarias (como el perro, con 4800 imágenes)
          ignorando por completo a las minoritarias (como el elefante, con 1400
          imágenes).

    \item El EarlyStopping con \texttt{patience=5} y
          \texttt{restore\_best\_weights=True} demostró ser una estrategia
          indispensable para prevenir el sobreajuste (overfitting).
\end{itemize}

\bibliographystyle{model1-num-names}
\bibliography{bibliografia}

\end{document}